% ============================================================================
%  ABSTRACT / SOMMARIO --- documento autonomo per la consegna (PDF/A)
%  Tesi di Laurea Magistrale --- Progetto HYDRAS
% ============================================================================
\documentclass[12pt,a4paper]{article}

\usepackage[utf8]{inputenc}
\usepackage[T1]{fontenc}
\usepackage[italian,english]{babel}
\usepackage{lmodern}
\usepackage[top=2.8cm, bottom=2.8cm, left=2.6cm, right=2.6cm]{geometry}
\usepackage{setspace}
\setstretch{1.2}
\usepackage[colorlinks=true, linkcolor=black, citecolor=black, urlcolor=black]{hyperref}
\hypersetup{
  pdftitle={Abstract --- Sviluppo di un sistema basato su reinforcement learning per la localizzazione di sorgenti inquinanti in ambiente marino},
  pdfauthor={Mattia Manneschi}
}

\pagestyle{empty}

\begin{document}

\begin{center}
    {\normalsize\scshape Universit\`a degli Studi di Firenze}\\[1pt]
    {\small\scshape Scuola di Ingegneria}\\[2pt]
    {\small\scshape Dipartimento di Ingegneria dell'Informazione}\\[1pt]
    {\small\scshape Dipartimento di Ingegneria Civile e Ambientale}\\[6pt]
    {\small Tesi di Laurea Magistrale in Ingegneria Informatica}\\[10pt]
    {\large\bfseries Sviluppo di un sistema basato su reinforcement learning per la
    localizzazione di sorgenti inquinanti in ambiente marino}\\[8pt]
    {\small\itshape Candidato:} Mattia Manneschi \qquad
    {\small\itshape Relatori:} Prof.\ G.\ Battistelli, Prof.\ N.\ Forti, Prof.ssa I.\ Simonetti\\[2pt]
    {\small Anno Accademico 2025/2026}
\end{center}

\vspace{4mm}
\hrule
\vspace{6mm}

\selectlanguage{italian}
\subsection*{Sommario}

L'inquinamento delle acque costiere richiede strumenti rapidi ed economici per individuare
la sorgente di uno sversamento. Questa tesi affronta il problema del \emph{source seeking}
marino --- la localizzazione autonoma della sorgente di un \emph{plume} inquinante da parte
di un veicolo subacqueo (AUV) --- confrontando due paradigmi di navigazione sullo stesso
ambiente di simulazione, costruito a partire da campi di concentrazione, vento e corrente
prodotti da simulazioni idrodinamiche ad alta risoluzione (MIKE21) di un tratto di mare
costiero.

Il primo approccio \`e il \emph{Field Climbing Method} (FCM), un metodo analitico che stima
il gradiente locale di concentrazione da una formazione di sensori e vi risale, nella sua
versione adattiva basata su Adam. Il secondo \`e un agente di \emph{Reinforcement Learning}
addestrato con PPO (nella variante MaskablePPO), che apprende una politica di navigazione
dalle sole osservazioni dei sensori. L'ambiente condiviso adotta una formazione
master--slave con una o due corone di sensori, uno spazio d'azione discreto con
\textbf{velocit\`a variabile appresa} e \emph{action masking} per i vincoli di costa e di
dominio.

I due metodi sono valutati su un insieme di sorgenti \emph{held-out}, mai viste in
addestramento, attraverso diversi scenari di vento/corrente e finestre temporali di
diffusione del plume (fino a 1560 episodi per modello). L'agente PPO raggiunge un tasso di successo
prossimo al \textbf{99\%} e supera nettamente e sistematicamente l'FCM (che non oltrepassa
il \textbf{64\%}), il cui limite si rivela \emph{strutturale} --- dovuto alla natura
reattiva e locale del metodo --- e non cinematico. La seconda corona di sensori si rivela un
vantaggio netto a partire dalle velocit\`a intermedie (da 2~m/s in su): riduce i tempi al
successo del 25--34\% a parit\`a di affidabilit\`a, mentre alle velocit\`a pi\`u basse
comporta un piccolo costo di robustezza.

Infine, la politica appresa \`e resa interpretabile tramite \emph{ablation} causale di
gruppo, valori di Shapley esatti e sonde comportamentali; un esperimento di riaddestramento
senza formazione ne quantifica il contributo. Il lavoro fornisce una pipeline completa dai
dati idrodinamici alla politica addestrata e interpretata, e una prova quantitativa del
vantaggio dell'apprendimento per rinforzo sui metodi reattivi di gradiente.

\medskip
\noindent\textbf{Parole chiave:} source seeking marino, veicoli autonomi subacquei,
reinforcement learning, PPO, metodo del gradiente, explainability, inquinamento marino.

\vspace{3mm}
\selectlanguage{english}
\subsection*{Abstract}

Coastal water pollution calls for fast and inexpensive tools to locate the source of a
spill. This thesis addresses marine \emph{source seeking} --- the autonomous localization of
the source of a pollutant plume by an autonomous underwater vehicle (AUV) --- by comparing
two navigation paradigms on the same simulation environment, built from concentration, wind
and current fields produced by high-resolution hydrodynamic simulations (MIKE21) of a
coastal sea area.

The first approach is the \emph{Field Climbing Method} (FCM), an analytical method that
estimates the local concentration gradient from a sensor formation and climbs it, in its
Adam-based adaptive variant. The second is a \emph{Reinforcement Learning} agent trained
with PPO (in the MaskablePPO variant), which learns a navigation policy from sensor
observations alone. The shared environment uses a master--slave formation with one or two
sensor rings, a discrete action space with \emph{learned variable velocity}, and action
masking for coast and domain constraints.

Both methods are evaluated on a \emph{held-out} set of sources, never seen during training,
across several wind/current scenarios and plume-diffusion time windows (up to 1560 episodes
per model). The PPO agent attains a success rate close to \textbf{99\%} and systematically
outperforms FCM (which never exceeds \textbf{64\%}), whose limitation proves to be
\emph{structural} rather than kinematic. The second sensor ring becomes a net advantage from
intermediate velocities (2~m/s and above): it shortens the time to success by 25--34\% at no
reliability cost, whereas at lower velocities it incurs a small robustness cost.

The learned policy is finally
made interpretable through causal group ablation, exact Shapley values and behavioral
probes, and a formation-free retraining experiment quantifies the formation's contribution.
The work delivers a complete pipeline from hydrodynamic data to a trained and interpreted
policy, and quantitative evidence of the advantage of reinforcement learning over reactive
gradient methods.

\medskip
\noindent\textbf{Keywords:} marine source seeking, autonomous underwater vehicles,
reinforcement learning, PPO, gradient method, explainability, marine pollution.

\end{document}
